\section{Problem Formulation}
\label{sec:problem}

\subsection{Setting: ambient contamination of a decision channel}
Consider an agent $A$ that, given a user instruction $u$, must choose one product
from a candidate list $\mathcal{C}=\{c_1,\dots,c_N\}$ rendered by a search or
recommendation layer. Each candidate is a \emph{listing}: a title, marketing
bullets, specification attributes, a description, a price, and an aggregate
rating with a review count. The agent observes a rendered version of these
fields and commits to a choice $k\in\{1..N\}$.

A note on scope: \emph{ambient} refers to how contamination arrives ---
synthetic content occupies the retrieval context that routine browsing
presents, with no adversary and no targeting --- not to a claim that the
Web as a whole is saturated. Our primary surface is the ten-slot results
page a search engine returns; a dedicated scale experiment extends the
surface to fifty slots and measures how degradation scales with the
absolute number of contaminated slots.

We model \emph{ambient synthetic contamination} as a random process on the
candidate pool, not as an adversary. A fraction $\varepsilon$ of distractor
slots is occupied by synthetic listings $s$ produced by a generative model
under \emph{commercial} optimization pressure: maximize apparent fit to the
user's intent (query keywords, bullet structure, rating and review volume),
minimize effort. Formally, for a task with gold product $g$ and real
distractors $D=\{d_1,\dots,d_{N-1}\}$, the contaminated pool is
$\mathcal{C}_\varepsilon = \{g\} \cup D' \cup S$, where $S$ replaces
$\lfloor\varepsilon (N{-}1)\rfloor$ elements of $D$. The gold product is never
replaced and the positions of all candidates are fixed per task, so that
measurements at different $\varepsilon$ are \emph{paired}.

Three properties distinguish this threat model from adversarial injection.
(i)~\emph{No targeting}: synthetic listings are produced for the category and
query in general, not against a specific user, task, or gold product; their
accidental overlap with any given instruction is incidental. (ii)~\emph{No user awareness}: generators see the category, the query, and
(for clones) a popular listing, but never the individual user's instruction,
so a mismatch with any particular user's required attributes is incidental
rather than engineered.
(iii)~\emph{Persistence}: such content is the by-product of ordinary incentives
(SEO, affiliate monetization), so its prevalence $\varepsilon$ is a property of
the Web, not of an attack session.

\subsection{Quantities of interest}
Let $r(c; u)\in[0,1]$ be the official reward of candidate $c$ under
instruction $u$, and let $K_\varepsilon\in\{1..N\}$ be the agent's choice at
contamination level $\varepsilon$. We measure:

\begin{itemize}
\item \textbf{Task success} $\mathrm{S}(\varepsilon)=\Pr[\,r(c_{K_\varepsilon};u)\ge
      1-\delta\,]$ with $\delta=10^{-3}$ (and a partial-success variant at
      $\delta=0.2$), and \textbf{regret}
      $\mathrm{R}(\varepsilon)=\mathbb{E}[\,1-r(c_{K_\varepsilon};u)\,]$, the
      economic loss of acting on the contaminated Web relative to the best
      available candidate.
\item \textbf{Susceptibility} $\sigma$, from a forced-choice probe: with only
      two candidates, the gold $g$ and one synthetic listing $s$, presented in
      randomized order, $\sigma=\Pr[K=s]$; the null value is
      $\sigma=\tfrac12$, and $\sigma>\tfrac12$ indicates a systematic
      preference for machine-optimized pages independent of any dose--response
      accumulation.
\item \textbf{Dose--response curve} $\mathrm{S}(\varepsilon)$ over
      $\varepsilon\in\{0, 0.25, 0.5, 0.75\}$, sufficient to locate the
      saturation point of the curve (the nominal $\varepsilon{=}0.1$ level
      maps to zero replaced slots and coincides with $0$).
\item \textbf{Calibration drift}: expected calibration error (ECE) of the
      agent's post-hoc choice confidence as a function of $\varepsilon$,
      capturing whether contamination distorts \emph{confidence} even as
      choices degrade.
\end{itemize}

\subsection{The silent-erosion hypothesis}
The security literature on web agents predicts degradation only when
adversarially targeted content is present. Ambient contamination implies a
different, more insidious signature: because synthetic listings are optimized
for \emph{perceived} quality rather than goal satisfaction, we expect the
choice distribution to shift toward them ($\sigma>1/2$; $\mathrm{R}$
increasing in $\varepsilon$) while coarse success rates may appear stable ---
regret rises through partial-quality purchases (attribute mismatches,
overpayment) long before complete failures appear. A dose--response with this
shape would mean the Web can lose decision quality at a nearly invisible rate
to the operators of agents, motivating provenance-ordered defenses rather
than detection alone. Our experiments are designed to falsify or confirm this
signature (Section~\ref{sec:results}).
